\section{Preliminaries}\label{sec:prelim}

\subsection{Quantum states and operations}

For a finite set $\Sigma$ of $n$ qubits, $\Hil_\Sigma = (\mathbb{C}^2)^{\otimes n}$.

\begin{definition}[Density operator]\label{def:density}
A \emph{partial density operator} on $\Hil_\Sigma$ is a linear operator $\rho$
on $\Hil_\Sigma$ that is positive, i.e.\ $\langle\psi|\rho|\psi\rangle \ge 0$
for all $|\psi\rangle \in \Hil_\Sigma$, and such that $\tr(\rho) \le 1$. It is a
\emph{density operator} when $\tr(\rho) = 1$. A pure state $|\psi\rangle$ is the
density operator $|\psi\rangle\langle\psi|$.
\end{definition}

\begin{definition}[Super-operator]\label{def:sop}
A \emph{super-operator} on $n$ qubits, written $\sop{E} : \mathsf{Sop}(n)$, is
a quantum channel: a completely positive and trace-preserving linear map on
the operators of $(\mathbb{C}^2)^{\otimes n}$. Equivalently,
$\sop{E}(\rho) = \sum_i E_i \rho E_i^\dagger$ for a family $\{E_i\}_i$ of
linear operators such that $\sum_i E_i^\dagger E_i = I$. For a tuple
$\tilde q$ of $n$ qubits of $\Sigma$, $\sop{E}^{\tilde q}$ is the super-operator
on $\Hil_\Sigma$ that applies $\sop{E}$ to $\tilde q$ and the identity to the
other qubits. Applying $\sop{E}$ to $\tilde q$ maps $\rho$ to $\sapp{\sop{E}}$,
and $\tr(\sapp{\sop{E}}) = \tr(\rho)$.
\end{definition}

The programs use two kinds of super-operators. A \emph{quantum gate} on $n$
qubits, written $U : \mathsf{Op}(n)$, is a unitary operator on
$(\mathbb{C}^2)^{\otimes n}$, i.e.\ $UU^\dagger = U^\dagger U = I$; it gives
the super-operator $\sop{E}_U : \rho \mapsto U \rho U^\dagger$, so that
$\sapp{\sop{E}_U} = \upd{U}$, where $U^{\tilde q}$ applies $U$ to $\tilde q$
and the identity to the other qubits. The \emph{initialisation} resets $n$
qubits to $|0\rangle$; it is the super-operator
$\sop{E}_{\mathsf{init}} : \rho \mapsto \sum_{i} |0\rangle\langle i| \,\rho\,
|i\rangle\langle 0|$, where $i$ ranges over the computational basis, and we
write $\inirho$ for $\sapp{\sop{E}_{\mathsf{init}}}$.

\begin{definition}[Measurement]\label{def:meas}
A \emph{measurement} on $n$ qubits with $k$ outcomes, written
$\M : \mathsf{Meas}(n)$ with $|\M| = k$, is a family $\{M_m\}_{0 \le m < k}$ of
linear operators on $(\mathbb{C}^2)^{\otimes n}$ such that
$\sum_{m} M_m^\dagger M_m = I$. Measuring $\tilde q$ in $\rho$ gives the outcome
$m$ with probability $p_m = \tr(\br)$, and the \emph{unnormalised branch}
$\br$, whose trace is $p_m$. The normalised post-measurement state is the branch
divided by $p_m$.
\end{definition}

\subsection{Syntax of lqCCS}

Let $\mathit{Tag} \ni t$ be a countably infinite set of \emph{tags}.

Outputs are asynchronous, sums are guarded, and recursion is as in VACCS, with
process variables $X$ and a tagged recursion operator:
\[
\begin{array}{r@{\ ::=\ }l}
P & \tagged{t}{c!e} \mid K \mid P \parallel P \mid \ite{e}{P}{P} \mid P \setminus c
    \mid X \mid \tagged{t}{\rec{X}{P}}\\
K & \mathbf{0}_{\tilde e} \mid \alpha.P \mid K + K\\
\alpha & \tagged{t}{\tau} \mid \tagged{t}{c?x} \mid
        \tagged{t}{U(\tilde e)} \mid \tagged{t}{\M(\tilde e \rhd x)}
        \mid \tagged{t}{\mathsf{init}(\tilde e)}
\end{array}
\]
The prefix $\tagged{t}{U(\tilde e)}$ applies the quantum gate $U$ to the
qubits $\tilde e$, and the prefix $\tagged{t}{\mathsf{init}(\tilde e)}$ resets
the qubits $\tilde e$ to $|0\rangle$. Both are super-operators
(\cref{def:sop}), and the memory treats them alike.
$P[Q/X]$ is the substitution of $Q$ for the free occurrences of $X$ in $P$;
unfolding a recursion is a step of its own, selected by the tag of the recursion.
A \emph{choice} is either a tag or an unordered pair of tags,
$s ::= t \mid \{t,t'\}$: the pair $\{t,t'\}$ selects the synchronisation
between a move tagged $t$ and a move tagged $t'$, whichever side each move is
on. A \emph{configuration} $\conf{\rho}{P}$ is a program $P$ together with
the quantum memory $\rho$ it acts upon.
